We reassess periodicity-hub organization in the R{\"o}ssler flow using
parameter-plane atlases, corrected periodic orbits, independent integrators,
Floquet analysis, and reconstructed chaotic saddles. Eleven sampled
$(a,c)$ planes at $b=0.10,0.12,\ldots,0.30$ reveal coherent changes in
shrimp-like periodic windows; these finite-time images provide context rather
than a global bifurcation classification.

At the reported Jones hub, the small equilibrium has the stated saddle-focus
eigenstructure. Its analytically determined Hopf locus anchors a continued
period-1 family, with four numerically qualified supercritical doublings
through a stable period-16 child on a declared fixed-$a$ path. A separate
manifold boundary-value calculation supports a nearby homoclinic candidate at
$(a,b,c)\simeq(0.182643608174,0.2,10.3084)$, with agreement across two
integrators, shooting meshes, and matching radii. This differs from the
printed coordinate; it neither excludes another connection there nor proves
homoclinic existence. Subsequent continuation exhibits a sampled parameter
turn, whose physical location still requires conditioning-aware error control.

On the Barrio--Blesa--Serrano section, independent saddle reconstructions
localize a change between two and three branches to
$a\in[0.1481875,0.14825]$ at $(b,c)=(0.2,20)$. Eleven recovered unstable
periodic-orbit families persist across this interval. A pre-existing unstable
lobe appears in the reconstructed saddle only on the three-branch side,
supporting a pruning or reinjection hypothesis that remains to be tested by
direct manifold geometry.

Continuation of a period-6 family resolves two arms of its flip locus and
distinguishes section grazing from an invariant orbit bifurcation. Its
returning family includes seven numerically qualified supercritical doublings
and evidence for a subcritical eighth birth. High-precision correction also
retracts an apparent daughter that collapses to a doubled parent. These
results strengthen local computational structure while leaving finite
logistic ordering, a continuous topology-transition curve, and a complete
explanation of the parameter plane open.
